\chapter*{Abstrak}

Ledakan informasi dalam platform hiburan digital menuntut sistem pemberi rekomendasi yang mampu memprediksi preferensi pengguna secara akurat di tengah tingginya tingkat kelangkaan data (\textit{data sparsity}). Penelitian ini menyajikan evaluasi komparatif komprehensif terhadap algoritma perekomendasi berbasis \textit{collaborative filtering} pada dataset MovieLens 100K (943 pengguna, 1.682 film, 100.000 rating, sparsitas 93,70\%) dengan skema protokol baku 5-\textit{fold cross-validation}. Seluruh algoritma dan metrik diimplementasikan secara modular dari awal (\textit{from scratch}) tanpa pustaka rekomendasi pihak ketiga, mencakup empat model \textit{baseline} (tebak acak, rata-rata global, rata-rata pengguna, dan rata-rata film), dua metode \textit{memory-based} (User-Based Collaborative Filtering dengan korelasi Pearson dan Item-Based Collaborative Filtering dengan \textit{adjusted cosine}), serta metode \textit{model-based} Matrix Factorization berbasis \textit{Stochastic Gradient Descent} (SGD) dengan regularisasi $L_2$. Evaluasi diukur melalui metrik kesalahan rating (MAE, MSE, RMSE, NMSE) dan kemampuan perankingan (NDCG@10). Hasil eksperimen menunjukkan bahwa \textit{baseline} rata-rata film merupakan acuan terkuat (MAE 0,8174; RMSE 1,0246; NDCG@10 0,8329). Metode User-Based CF menghasilkan MAE terendah (0,7342, perbaikan 10,18\% dibanding \textit{baseline}). Sementara itu, Matrix Factorization unggul secara konsisten dan signifikan pada seluruh lima \textit{fold} untuk metrik RMSE (0,9249, perbaikan 9,73\%, $p = 0,0002$ pada uji t berpasangan) serta kualitas perankingan daftar rekomendasi NDCG@10 (0,8578, perbaikan 2,99\%, $p = 0,0102$). Selain itu, Matrix Factorization menawarkan latensi prediksi tercepat (0,014 detik), menjadikannya kandidat paling optimal untuk deployment sistem rekomendasi skala besar.

\vspace{0.5 cm}
\begin{flushleft}
{\textbf{Kata Kunci:} Sistem Pemberi Rekomendasi, Collaborative Filtering, Matrix Factorization, MovieLens 100K, Evaluasi 5-Fold Cross Validation.}
\end{flushleft}
